\documentclass{ctexart}
\usepackage{Note}
\usepackage[
    backend=biber,
    style=chem-acs,
]{biblatex}
\addbibresource{ref.bib}
\usepackage{unicode-math}
\setCJKmainfont[
    BoldFont = 思源宋体 Bold,
    ItalicFont = 楷体,
]{宋体}
\author{邵天诚\ 2400011710\and 蒋锦豪\ 2400011785}
\title{\tbf{冰的微观结构}}
\date{}
\begin{document}\pagestyle{empty}
\maketitle
\section{晶态冰中的无序结构\cite{Salzmann2019}}
\subsection{氢原子的无序性}
在部分晶态冰中,氧原子形成规则的晶格结构,但氢原子的位置却存在无序性.以典型的冰$\text{I}_\text{h}$为例.考虑到六方冰中O原子做六方密堆积排列,因此O的位置就固定不动.由于固态的冰中存在氢键网络,每个O原子都通过氢键和周围四个O原子连接,因此我们需要做一个假设,即每个O原子周围都有两个H与其距离较远,另外两个与其距离较近.这对应着O形成两根O$-$H化学键和两根O$\cdot\cdot$H氢键.\\
\indent 这样,我们只需要考虑H的位置即可.对于$\SI{1}{\mole}$冰中的$\SI{2}{\mole}$\ce{H}原子,都有两种状态,即处于两个O原子之间离其中某个O原子更近的位置.这样的微观状态数一共有$2^{2N_{\text{A}}}$种.考虑到O原子对H原子的位置,每个O原子周围恰好有两个H靠近,两个H远离.这样的概率为
\[P_O=\dfrac{C_4^2}{2^4}=\dfrac{3}{8}\]
于是总的微观状态数为
\[\Omega=2^{2N_{\text{A}}}\cdot\left(P_O\right)^{N_{\text{A}}}=\left(\dfrac32\right)^{N_{\text{A}}}\]
于是
\[S=k_\text B\ln\Omega=R\ln\dfrac32\si{\joule\per\kelvin}\approx3.37\si{\joule\per\kelvin}\]
实际上可以做改进\cite{Li2024}.
\indent 冰$\text{I}_{\text{h}}$并非唯一的氢无序相,所有由液态水结晶形成的冰相均为氢无序.将氢无序冰冷却,可观测到其转变为对应的氢有序相.目前仅有冰 III与冰VII在冷却时可自发完成氢有序化,而其余有序冰的制备需借助酸碱掺杂.在氢有序化过程中,氢键网络拓扑结构保持不变.冰II是唯一无氢无序对应相的氢有序冰相:理论计算预测,冰II与其氢无序对应相冰$\text{II}_\text{d}$的相界位于液态水稳定存在的区域.压缩冰 VII至冰X时,氢键变为对称氢键,氢有序-无序现象不再适用。
\subsection{冰I型的无序性}
所谓立方冰$\text{I}_\text{c}$样品并非纯立方结构,同时包含六方结构.这类堆垛无序冰I(称作$\text{I}_{\text{sd}}$)存在两种堆垛方式,使相邻层间的六元环呈现不同构象\\
\indent 近十年,该领域研究因堆垛无序漫衍射特征模拟技术的进步取得重大突破.基于相关分析,可测定冰$\text{I}_{\text{sd}}$的立方度(即立方堆积占比).不同制备方法可得到立方度差异极大的样品.常压下加热冰II制得的冰$\text{I}_{\text{sd}}$立方度达$73.3\%$;真空下过冷水纳米液滴冻结得到的冰$\text{I}_{\text{sd}}$立方度为$78\%$,是目前已知立方度最高的样品.\\
\indent 除冰$\text{I}_\text{c}$外,还可探索具有高阶记忆效应的冰I型多型体(立方与六方堆垛严格按一定规律交替).
\subsection{掺杂导致的有序/无序化}
向水中掺杂~\ce{HCl},~\ce{LiOH}等酸或碱可以引导水分子重定向,形成有序度更高的冰晶体.右图以焓变示意了氢无序冰冷却时的不同转变路径.高温下,氢无序冰中水分子发生动态重排;冷却至$T_g$时,部分冰的重排动力学显著减慢,偏离平衡态形成玻璃态,对应路径(1).向冰种掺杂可以引入外源性的~\ce{H3O+}缺陷加速重排,得到部分氢有序的冰(路径(3))或完全有序的冰(路径(4)).\\
\indent 研究提出决定掺杂剂效果的三大因素:掺杂剂需在冰中高浓度溶解;在冰中具备强酸性或碱性,产生可加速重排的~\ce{H3O+}或~\ce{OH-}缺陷;~\ce{H3O+}或~\ce{OH-}缺陷的迁移动力学可能存在本质差异.\\
\indent \ce{NH4+}与~\ce{F-}掺入氢有序冰(如冰 II)会引发氢无序:在无缺陷,氢有序的冰II中,离子充当 ``拓扑电荷'',长程破坏水分子的取向有序,导致冰II自由能大幅升高,这也是掺杂一定量~\ce{NH4F}时冰II从相图中消失的原因.
\section{无定形冰\cite{Loerting2011}}
无定形冰与晶态冰存在显著差异,核心在于无定形冰缺乏长程有序性,没有平移,旋转等对称性.
\subsection{无定形冰的相图}
此处的\textit{相图}不是真正意义上的相图,其中的状态不是热力学最稳定的状态,而只是动力学上能够存在的状态.从上至下依次为液态水,过冷水,\,``无人区'',超粘性态与无定形态.\,``无人区''指在该条件下极易形成晶体,难以形成无定形态.超粘性态是无定形态升温至高于玻璃化温度得到的状态,极其不稳定.后面将主要讨论无定形态.
\subsection{无定形冰的合成}
无定形冰的合成主要有以下几种方法.
\begin{enumerate}
    \item \tbf{气相沉积法}\\
    在超高真空(UHV)环境下,将水蒸气以极慢的速度喷射并沉积到温度极低$(<\SI{77}{\kelvin})$的金属基底上.
    \item \tbf{液体超快淬冷法}\\
    将液态水雾化成极其微小的纳米级或微米级水滴,再将这些微滴以极高的速度喷射到处于液氮温度的冷盘上.
    \item \tbf{高压诱导非晶化法}\\
    在低温$(<\SI{77}{\kelvin})$下,高压$(>\SI{0.1}{\giga\pascal})$压缩晶态冰,得到无定形冰.
\end{enumerate}
\subsection{无定形冰的分类}
历史上通过各种方法得到了许多不同的无定形冰并创造了许多名称,如何将这些无定形冰分类是一个核心问题.以下是分类的依据:
\begin{enumerate}
    \item \tbf{径向分布函数}\\
    径向分布函数:以某一个中心原子为原点,在距离为$r$处发现另一个原子的局域概率密度,与全局平均密度的比值.对于晶体而言,其径向分布函数为$\delta$函数(原子只在特定距离出现);对于无定形态,原子距离具有一定随机性,在径向分布函数中呈现出一定的峰宽.通过径向分布函数,可大致将无定形冰分为三类.
    \item \tbf{退火实验(核心判据)}\\
    非晶态结构会自然发生弛豫(分子间进行小范围重排)使能量降低,升高温度可加快该过程,通过升温使无定形态结构达到``平衡态''的过程称为``退火''.此处的``平衡态''依然不是热力学最稳定的状态,热力学最稳定的状态仍然是晶态,退火过程需要小心操作避免结晶.在不同压力下进行退火操作可得到该压力下最稳定的无定形态,以该状态下的密度对压力作图,可明确将无定形态划分为三种不同的状态,分别为LDA (Low-density amorphous ice)、HDA (High-density amorphous ice)、VHDA (Very high-density amorphous ice).这是进行状态划分的核心依据.
\end{enumerate}
\subsection{不同无定形态之间的转化}
多种不同的实验证明在常压下VHDA与HDA会自发向LDA转化,高压下HDA会自发向VHDA转化.以其中的两个实验为例,左图所示实验为示差扫描量热法,峰代表在该温度下出现明显的放热；右图实验为拉曼光谱法,展现升温过程中~\ce{O-H}键振动波数的变化.
\subsection{无定形冰的势能曲线}
依据上述实验,作者构筑出三种不同压力下无定形冰的势能曲线.常压下VHDA不存在势能极小值点, $\SI{1.2}{\giga\pascal}$下LDA不存在势能极小值点,而$\SI{0.5}{\giga\pascal}$下三种无定形冰的势能极小值点同时存在.
\section*{分工}
蒋锦豪负责第一部分:晶态冰中的无序结构;邵天诚负责第二部分:无定形冰.
\printbibliography[title={\Large 参考文献}]
\end{document}